This research builds upon existing applications of AI, IoT, and remote sensing in Indian agriculture. We draw from Vaswani et al. (2017) for transformer architectures and Jain et al. \cite{jain2022iot} for IoT-based soil monitoring, while leveraging ISRO's Bhuvan platform and Ministry of Agriculture data \cite{agriculture2023ministry} for real-world crop assessment.

Remote sensing platforms like Bhuvan enable large-area crop monitoring through high-resolution temporal imagery, facilitating early stress detection and yield estimation as demonstrated by Sharma et al. Recent AI advances show transformer models effectively combine multi-modal data; Li et al. (2025) and Mehta et al. \cite{mehta2024multimodal} prove their value in analyzing NDVI and thermal indices for major crops. The foundational Video Swin Transformer architecture \cite{liu2022swin3d}, building upon the original Swin Transformer \cite{liu2021swin}, enables efficient spatiotemporal feature extraction through shifted window self-attention — a key component of our temporal processing pipeline.

Federated learning addresses India's fragmented land holdings while preserving privacy. Kumar et al. \cite{kumar2023federated} show its utility for yield prediction without data transfer, while Das and Borthakur \cite{das2023federated} demonstrate its applications in crop disease detection using localized datasets across Indian regions. Edge computing enables rural accessibility through TinyML implementations like Reddy et al.'s \cite{reddy2023tinyml} animal intrusion detection and analysis of low-power agricultural IoT.

Graph Neural Networks have emerged as a powerful tool for modelling spatial farm dependencies. Hamilton et al. \cite{hamilton2017graphsage} introduced GraphSAGE enabling inductive learning on unseen graph nodes, which is critical in agricultural contexts where new farm plots are continuously registered. Brody et al. \cite{brody2022gatv2} further improved expressiveness through GATv2's dynamic attention mechanism, which allows attention coefficients to vary based on both source and destination node features — essential for capturing heterogeneous soil and crop conditions across neighbouring fields.

Model compression for edge deployment is validated by Han et al. \cite{han2016deep}, who established pruning and quantization pipelines reducing model size by up to $4\times$ with less than $1\%$ accuracy degradation. Qin et al. \cite{qin2023onnxquant} further validated INT8 ONNX dynamic quantization for real-time inference on ARM-based mobile devices, demonstrating its suitability for deployment on mid-range Android hardware — the same target platform used in this project.

Finally, data quality and completeness directly affect model reliability. Patel et al. \cite{patel2024imputation} validate predictive mean matching for handling gaps in pest monitoring records, while Krishnan et al. \cite{krishnan2023reconstruction} address sensor malfunction correction in Tamil Nadu farm datasets using deep learning reconstruction. These works highlight the importance of imputation strategies when working with real-world agricultural data; our system uses temporal median imputation to fill cloud-cover gaps in satellite imagery, which follows a similar rationale.
